\section{Inference and Deployment}

\subsection{Inference Pipeline}

Describe how a trained model or finalized method is used at test time.
State the required inputs, preprocessing steps, decoding or post-processing steps, and expected outputs.
If the inference procedure differs from training, explain the difference clearly.

\begin{algorithm}[t]
  \caption{\textbf{Illustrative inference procedure.} Replace the placeholder steps with the method's actual algorithm.}
  \label{alg:inference}
  \begin{algorithmic}[1]
    \Require Input sample $x$, trained model $f_{\theta}$
    \Ensure Prediction $\hat{y}$
    \State $z \gets \textsc{Preprocess}(x)$
    \State $s \gets f_{\theta}(z)$
    \If{$\textsc{Confidence}(s) < \tau$}
      \State $s \gets \textsc{Refine}(z, s)$
    \EndIf
    \State $\hat{y} \gets \textsc{Decode}(s)$
    \State \Return $\hat{y}$
  \end{algorithmic}
\end{algorithm}

\subsection{Efficiency}

Report runtime, memory usage, model size, or computational complexity when relevant.
Use a table if several variants are compared.
Avoid hard-coded claims until the final benchmark numbers have been verified.

\subsection{Practical Considerations}

Mention constraints that matter for real use, such as input resolution, supported formats, batching behavior, failure cases, or compatibility with downstream tools.
